\section{Day 14: Primitive Roots of Unity (Nov.\ 3, 2026)}
Last class, we started analyzing $K = K_{p^a}$, where $p$ is prime. We have that
\[ \abs{K:\QQ} = \varphi(p^a) = p^{a-1}(p - 1), \]
and that $\Gal(K/\QQ) = (\ZZ/p^a)^\times$. We also saw that
\[ \Phi_{p^a}(x) = \frac{x^{p^a} - 1}{x^{p^{a-1}} - 1} \]
is irreducible. Now, define $\lambda = 1 - \zeta_{p^a}$; since
\[ N(\lambda) = \prod_{b \in (\ZZ/p)^\times} (1 - \zeta^b) = \Phi_{p^a}(1) = p, \tag{$\ast$} \]
it follows that $(\lambda)$ is a prime ideal. Note that
\[ \left(\prod_{b \in (\ZZ/p)^\times} \frac{1 - \zeta^b}{1 - \zeta}\right)(1 - \zeta)^{\varphi(p^a)} = p \]
by ($\ast$), and so $\smash{\frac{1 - \zeta^b}{1 - \zeta} = 1 + \zeta + \dots + \zeta^{b-1} \in \SO_K}$ is a unit. We now want to show that $\smash{\frac{1 - \zeta}{1 - \zeta^b} \in \SO_K}$. First, $p \nmid b$, so there exists $c$ such that $bc = 1$ modulo $p^a$, and so $\zeta = (\zeta^b)^c$. We also have that
\[ \frac{1 - \zeta}{1 - \zeta^b} = \frac{1 - (\zeta^b)^c}{1 - \zeta^b} = 1 + \zeta^b + \dots + (\zeta^b)^{c-1} \in \SO_K, \]
whence $p = u \lambda^{\varphi(p^a)}$, where $u \in \SO_K^\times$. Writing $d = \varphi(p^a)$, we have that $p \SO_K = \lambda^d \SO_K$, where $(\lambda)$ was seen to be prime as before, so $e((\lambda) \mid p) = d$ and $f((\lambda) \mid p) = 1$. Thus, $\SO_K / (\lambda) \cong \FF_p$. Last time, we saw that
\[ \abs{\Delta(1, \zeta, \dots, \zeta^{d-1})} = p^{p^{a-1}(a (p-1) - 1)} = p^{a \varphi(a^p) - \frac{\varphi(a^p)}{p-1}} = p^{ad - \frac{d}{p-1}}. \]
In particular, all other primes are unramified in $K$.
\begin{theorem} \label{thm:cyclotomic_number_field_is_extension_by_root_of_unity}
    $\SO_K = \ZZ[\zeta]$.
\end{theorem}
\begin{proof}
    Since $\SO_K / (\lambda) \cong \FF_p$, for $x \in \SO_K$, we have that $x = c_0 + \lambda x_1$, with $x_1 \in \SO_K$ and $c_0 \in \ZZ$ with $0 \leq c_0 < p$. Repeating this $d$ times, we get
    \[ x = c_0 + c_1 \lambda + \dots + c_{d-1} \lambda^{d-1} + \lambda^d x_d, \]
    where $x_d \in \SO_K$, and the first $d$ terms lie in $R \coloneq \ZZ[\zeta]$. Since $(\lambda^d) = (p)$, it follows that $\lambda^d x_d = p x_d'$, where $x_d' \in \SO_K$.

    The upshot is that $\SO_K = R + p \SO_K$. Now, let $T = \SO_K/R$ be an abelian group. $T$ is finite and its order is a power of $p$. Also, note that $p \colon T \to T$ is a bijection (where $p$ denotes multiplication by $p$), so for $\ol x \in T$, where $x = x_0 + p x_1$ with $x_0 \in R$, we have that $\ol x = p \ol x_1$. We get a contradiction out of this; if $\abs{T}$ is a power of $p$, then $p \colon T \to T$ is not injective, so $\abs{T} = 1$, i.e., $\SO_K = R$.
\end{proof}

We now discuss the general roots of $1$.
\begin{theorem} \label{thm:roots_of_1}
    For all $u$, we have that $\abs{K_n : \QQ} = \varphi(n) \eqcolon d$, where $\SO_{K_n} = \ZZ[\zeta_n]$, and
    \[ \abs{\Delta_{K_n}} = \frac{n^d}{\prod_{p \mid n} p^{\frac{d}{p-1}}}. \]
\end{theorem}
If $p \mid n$, then the power of $p$ in $\abs{\Delta_{K_n}}$ is $ad - \frac{d}{p-1}$, where $a$ is the largest choice such that $p^a \mid n$.
\begin{semiproof}
    We will induct on the number of distinct primes dividing $n$. For the base case of $1$ prime, we have $n = p^a$, which we've already resolved. Now, write $n = p^a m$, where $\gcd(p,m) = 1$ with $m > 1$. Let $E = K_{p^a}$ and $F = K_m$. Observe that $EF = K_n$; $K_n$ is generated by $\zeta_n$, which is a primitive $n$th root of unity. Since $p$ and $m$ are coprime, we have that $um + vp^a = 1$ for some $u, v \in \ZZ$, so $\zeta_{p^a}^u \zeta_m^v$ is a primitive $n$th root of unity as well, so $K_n \subset EF$.

    Conversely, $\zeta_n^m$ is a primitive $p^a$th root of unity, and $\zeta_u^{p^a}$ is a primitive $m$th root of unity. This means $E \subset K_n$ and $F \subset K_n$.
\end{semiproof}

Before we finish the proof, let us make some observations first (which are left as an exercise).
\begin{enumerate}[(i)]
    \item $E \cap F = \QQ$.
    \item $\abs{EF : \QQ} = \abs{E : \QQ} \abs{F : \QQ} = \varphi(p^a) \varphi(m) = \varphi(n)$. 
\end{enumerate}

\begin{proof}[Proof of \ref{thm:roots_of_1}, cont.]
    The upshot being that $\abs{K_n : \QQ} = \varphi(n)$. By induction, $\Delta_{K_m}$ and $\Delta_{K_{p^a}}$ are coprime. This yields
    \begin{enumerate}[(i)]
        \item $\SO_{K_n} = \SO_{K_m} \SO_{K_{p^a}} = \ZZ[\zeta_m, \zeta_{p^a}] = \ZZ[\zeta_n]$,
        \item $\Delta_{K_n} = \Delta_{K_m}^{\abs{K_{p^a} : \QQ}} \cdot \Delta_{K_{p^a}}^{\abs{K_m : \QQ}}$.
    \end{enumerate}
    The power of $p$ in $\Delta_{K_n}$ is given by
    \[ \Delta_{K_{p^a}} = a p^{\varphi(p^a) - \frac{\varphi(p)}{p-1}} = \left(a \varphi(p^a) - \frac{\varphi(p^a)}{p - 1}\right) \varphi(m) = a \varphi(n) - \frac{\varphi(n)}{p - 1} = ad - \frac{d}{p-1}. \qedhere \]
\end{proof}

\begin{corollary}
    For all $n$, the polynomial
    \[ \Phi_n(x) = \prod_{\substack{1 \leq b \leq n \\ \gcd(b, n) = 1}} (x - \zeta_n^b) \]
    is the minimal polynomial of $\zeta_n$. In particular, $\Phi_n(x) \in \ZZ[x]$ and is irreducible over $\QQ$. Moreover, $\Gal(K_n/\QQ) = (\ZZ/n)^\times = \sigma_b(\zeta_n) = \zeta_n^b$, and
    \[ x^n - 1 = \prod_{d \mid n} \Phi_d(x). \]
\end{corollary}
